\chapter{基础知识}

\section{模型预测控制}

模型预测控制是一类能够在满足系统约束条件的前提下求解最优控制输入
的先进控制方法，其核心特征体现在后退时域控制（Receding Horizon Control, RHC）的思想上\cite{review_mpc}。具体而
言，MPC 在每个采样时刻均基于被控对象的模型在有限预测时域内对系统状态进行滚动预测，并在此基础上构建
带约束的优化问题，以求得对未来一段时间内最优的控制序列。随后，控制器仅执行该最优控制序列中的第一个
控制量，并在下一采样时刻利用新的观测状态重新建立预测与优化过程，从而形成闭环控制策略。

\begin{figure}
    \centering
    \includegraphics[width=.75\linewidth]{survey/mpc.png}
    \caption{\label{fig:mpc}模型预测控制\cite{mpc_101}}
\end{figure}

如\autoref{fig:mpc}所示，在时刻$k$，参考轨迹以红线表示，当前观测状态以黄色标识。优化问题在区间$[k, k+p]$上建立，
其中$p$为预测时域长度。典型情况下，MPC 的性能指标可以使用参考轨迹与预测轨迹之间偏差以及控制输入变化量
构造二次型代价函数，如：

\begin{equation}
\label{equ:mpc_cost}
\begin{aligned}
    J &= \sum_{i=0}^{p-1} \left( \| x_{k+i} - x_{ref,k+i} \|_Q^2 
    + \| u_{k+i} - u_{ref,k+i} \|_R^2 \right) 
    + \| x_{k+p} - x_{ref,k+p} \|_S^2 \\
    &\quad + \sum_{i=0}^{p-2} \| u_{k+i+1} - u_{k+i} \|_T^2
\end{aligned}
\end{equation}

其中，第一项用于惩罚状态偏差和控制偏差，第二项用于惩罚终端状态误差，第三项则是惩罚控制变化。通过求解上述优化问题，
可得到未来时域的最优控制序列，如图中绿色所示。随后，控制器执行该序列中的第一个控制输入，并在下一时刻基
于新的系统状态重复上述预测和优化过程，由此形成不断前移的预测窗口，这正是 RHC 思想的体现。

随着计算能力的提升，MPC 已逐渐拓展至多领域的复杂系统控制，包括机器人、航天器以及智能交通系统等。在结构上，
MPC 可分为线性 MPC（LMPC）、非线性 MPC（NMPC）、鲁棒 MPC 以及分布式/分层 MPC 等不同类型，分别针对线性系统、
非线性系统、模型不确定性以及大规模系统优化问题提供有效解决方案。
在自动驾驶控制中，MPC 已得到广泛应用。例如，在轨迹跟踪场景中，MPC 可根据车辆动力学模型和参考轨迹实时计算最优控
制输入，同时考虑速度、转向和环境约束；在障碍物规避任务中，MPC 可通过引入状态约束确保车辆避开静态或动态障碍物；
在速度与能耗优化方面，MPC 可以综合考虑驾驶舒适性、安全性与能量消耗，实现多目标协调控制。通过这些应用，MPC 展现出
处理复杂约束、提升系统性能以及适应非线性动力学的显著优势，也成为自动驾驶控制领域的重要研究工具。

\section{高斯过程回归}

高斯过程回归（Gaussian Process Regression, GPR）是一种非参数的贝叶斯统计学习方法，用于对未知函数进行预测与不确定性量化。
与传统的参数化回归方法不同，GPR 不假设具体的函数形式，而是通过协方差函数（核函数）刻画输入空间
中不同点之间的相关性，从而在给定观测数据下构建函数分布的后验。

高斯过程完全由其均值函数和协方差函数来表示\cite{gpml}。
我们用$m(x)$和$k(x, x^{'})$表示一个过程 $f(x) \sim \mathcal{GP}(m(x), k(x, x^{'}))$ 的均值和协方
差：
\begin{equation}
\label{equ:gpr_intro1}
\begin{aligned}
    m(x) &= \mathbb{E} [f(x)] \\
    k(x, x^{'}) &= \mathbb{E} [(f(x) - m(x))(f(x^{'}) - m(x^{'}))] \\
\end{aligned}
\end{equation}

具体而言，设训练数据集为 
$\mathcal{D} = \{ (x_i, y_i) \}_{i=1}^N$，GPR 假设函数值 $f(x)$ 在每个输入点 $x$ 上服从高斯分布，
并通过协方差矩阵 $K$ 描述点间相关性，则对于新的输入 $x^*$，预测输出 $y^*$ 的分布为

\begin{equation}
\label{equ:gpr_intro2}
\begin{aligned}
    & f^* \mid X, y, x^* \sim \mathcal{N}(\mu^*, \Sigma^*) , \\
    & \mu^* = K_{*f} (K_{ff} + \sigma_n^2 I)^{-1} y, \\
    & \Sigma^* = K_{} - K_{*f} (K_{ff} + \sigma_n^2 I)^{-1} K_{f*} ,
\end{aligned}
\end{equation}

其中 $K_{f*}$ 是训练点与测试点之间的协方差向量，$K_{ff}$ 为训练点之间的协方差矩阵，$\sigma_n^2$ 为观测噪声方差。
通过该分布，GPR 不仅能够给出预测均值，还能够量化预测的不确定性，这对于涉及安全性和鲁棒性的自动驾驶控制任务尤为重要。

\begin{figure}
    \centering
    \includegraphics[width=.75\linewidth]{survey/gpr.png}
    \caption{\label{fig:gpr_intro}高斯过程回归在无噪声数据中的预测\cite{gpr_intro}}
\end{figure}

GPR 的优势在于：首先，它能够自然地处理小样本学习问题，在数据有限的情况下仍能提供有效预测；其次，贝叶斯框
架允许对预测不确定性进行量化，从而为控制策略提供风险评估信息；最后，通过选择不同的核函数，可以灵活建模不
同类型的非线性函数关系。由于这些特性，GPR 已被广泛应用于自动驾驶中的动力学建模、轨迹预测以及基于学习的控制
器设计等场景。例如，在基于模型的强化学习和非线性 MPC 中，GPR 可以用来在线更新系统模型，提升预测精度和控制性能。

\section{基于学习的模型预测控制}

随着机器学习技术在控制领域的快速发展，基于学习的模型预测控制（Learning-based MPC, LB-MPC）逐渐成为处理复杂系统、
不确定性建模以及高维控制任务的重要手段。与传统 MPC 严格依赖由物理机理推导的精确模型不同，LB-MPC 
将数据驱动方法融入 MPC 的建模与优化流程，通过学习补偿系统动态、自动调整优化项或在控制回路中引入自
适应机制，以提高控制器在未知环境与动态变化下的鲁棒性和性能表现\cite{survey_lb_mpc}。因此，LB-MPC 代表了 MPC 从“基于
模型”走向“模型与学习融合”的重要趋势。

总体而言，LB-MPC 的核心思想是在 MPC 框架中引入学习模块，以更好地描述系统真实动力学、优化控制策略或
增强安全性。其典型研究方向包括以下几个方面：

\begin{itemize}
    \item 学习系统模型: 传统 MPC 实现的关键在于获得系统的精确预测模型。然而对于高度非线性、强耦合或具有时变特性的复杂系统，
基于物理机理的建模往往存在简化误差。LB-MPC 通过数据驱动的方法（如高斯过程回归、深度神经网络、稀疏回归等）
对系统动力学进行学习或补偿，从而构建更贴近真实系统的预测模型。
这类方法既可以完全用学习模型替代原模型，也可采用“残差学习”（learning residual dynamics）方案，即在已
有模型基础上学习未建模动态，从而显著提升预测精度并减少模型偏差对 MPC 性能的影响。
    \item 学习控制策略与优化结构: 除了系统模型，学习机制也可用于 MPC 优化问题中的关键结构。例如，通过学习终端代价、终端
    约束、控制权重、松弛因子等，可以增强 MPC 的性能和计算效率。部分研究通过深度强化学习（Deep RL）近似 MPC 的求解器，或通过模仿
学习（Imitation Learning）训练一个“快速求解策略”来替代在线优化，从而显著降低计算开销，为实时控制提供可能。此
外，基于学习的调参方法使得控制框架能够根据运行环境自动调整控制增益或代价函数结构，提高适应性。
    \item 安全学习与在线适应: 在许多安全关键场景（如自动驾驶、机器人操作、航空航天等），直接使用学习模型可能引入不确定
    性或潜在风险。因此 LB-MPC 中常采用“安全过滤”（safety filter）、“屏障函数”（CBF）或“鲁棒 MPC 结构”等机制，在利用
学习模型改善性能的同时确保系统满足安全约束。
同时，在线学习（online learning）方法可利用运行过程中获得的实时反馈不断更新模型或策略，使控制器能够
适应环境变化和外部扰动。例如，基于贝叶斯更新的方法可动态调整模型后验分布，深度学习模型可通过经验回放
缓慢更新参数，而基于元学习（meta-learning）的方案则可在结构层面实现快速适应。
\end{itemize}

根据学习发生的时机，LB-MPC 可分为：
\begin{itemize}
    \item 离线学习（offline learning）：利用大量采样数据训练高精度模型或策略，再嵌入 MPC 框架中进行控制。优
点是可使用高算力训练复杂模型；缺点是泛化能力有限。

    \item 在线学习（online learning）：在实际运行中实时更新模型或策略，具有较强适应性，但对计算资源及安全保障提出更高要求。

    \item 混合学习（hybrid learning）：结合离线与在线学习，先获得基础模型，再通过在线补偿提升精度与适应性，是当前应用最广的方案。
\end{itemize}